\documentclass[11pt]{article}
\usepackage[margin=1in]{geometry}
\usepackage{amsmath,amsfonts,amssymb,amsthm,bm,mathtools}
\usepackage{graphicx,booktabs,array,longtable}
\usepackage[hypertexnames=false]{hyperref}
\usepackage{float}
\usepackage{enumitem}
\usepackage{physics}

\newtheorem{theorem}{Theorem}[section]
\newtheorem{proposition}[theorem]{Proposition}
\newtheorem{definition}[theorem]{Definition}
\newtheorem{corollary}[theorem]{Corollary}
\newtheorem{remark}[theorem]{Remark}

\title{Hybrid Quantum--Classical Machine Learning for Computer Vision:\\ A Mathematical Literature Review}
\author{Sparsho~Chakraborty\\[4pt]
\small School of Basic Sciences, Department of Physics, Indian Institute of Technology Bhubaneswar, Odisha, India.}
\date{}

\begin{document}
\maketitle

\begin{abstract}
This review surveys, with explicit mathematical formalism, the technical foundations underlying hybrid quantum--classical machine learning for computer vision. We cover: variational quantum algorithms and the parameter-shift rule for gradient estimation; the barren-plateau phenomenon, its variance-scaling derivation, and architectural, symmetry-restricted, and Hamiltonian-generated mitigation strategies; quantum representations of classical images (FRQI, NEQR) and their state-preparation gate complexity; quantum convolutional neural networks and their multiscale-entanglement-renormalization-derived pooling structure; tensor-network methods for machine learning, including the matrix-product-state (MPS) formalism, canonical form via singular value decomposition (SVD), entanglement entropy as a compression and feature-relevance measure, and Fourier-spectrum conditions for efficient MPS-to-circuit loading; quantum autoencoders, dressed (classical-pre/post-processed) quantum circuits, and quantum generative models, with their fidelity-, adversarial-, and diffusion-based training objectives; quantum kernel methods and the reproducing-kernel-Hilbert-space view of quantum-enhanced feature maps; symmetry and group-equivariance in classical and quantum vision models, developed from first principles through representation theory, with the dihedral group $D_4$ worked through explicitly, including its character table; physics-informed and Hamiltonian-based regularization, including the algebraic structure and Bethe-ansatz integrability of the anisotropic Heisenberg (XYZ) spin chain and its $\mathrm{SU}(2)/\mathrm{U}(1)$ symmetry hierarchy; and the quantum-advantage and dequantization literature that determines when any of the preceding architectural choices can be expected to yield a genuine advantage over a resource-matched classical alternative. Approximately seventy primary sources, each independently verified against its original venue rather than reproduced from a secondary summary, are synthesized into a single mathematically self-contained narrative, closing with a comparative summary table spanning the architecture families surveyed.
\end{abstract}

\tableofcontents

\section{Introduction}
\label{sec:intro}
Hybrid quantum--classical machine learning couples a parameterized quantum circuit (PQC), executed on near-term noisy intermediate-scale quantum (NISQ) hardware \cite{Preskill2018}, to a classical optimizer and, frequently, a classical pre- or post-processing network. For computer vision specifically, this hybridization has been explored along at least seven largely independent technical axes: (i) how a classical image is represented as, or mapped into, a quantum state or quantum-derived feature vector; (ii) how a variational circuit is trained given the barren-plateau trainability obstruction; (iii) how a convolution-like, hierarchically pooled circuit structure can be built directly, independent of the barren-plateau mitigations developed for generic ansätze; (iv) how tensor-network structure — a classical simulation tool for quantum many-body systems — can itself be repurposed as a machine learning primitive, and how it interfaces with the gate complexity of loading a compressed representation onto real qubits; (v) how quantum autoencoders, dressed classical-quantum architectures, generative adversarial networks, and diffusion models are constructed as quantum or hybrid analogues of their classical counterparts; (vi) how symmetry and group-equivariance, well studied in classical geometric deep learning, can be built into variational quantum circuits; and (vii) under what precise conditions, if any, a quantum or quantum-inspired model exhibits an advantage over a resource-matched classical model, given the dequantization literature that has repeatedly narrowed the space of such claims.

This review treats each axis with explicit mathematical formalism rather than as a citation list, because the central methodological difficulty in this literature — as several of the sources below explicitly discuss \cite{Huang2021PowerData,Cotler2021,Gyurik2023} — is that architectural motivation and demonstrated advantage are frequently conflated. A circuit can be well-motivated by symmetry, physical analogy, hierarchical structure, or expressivity while still not being \emph{necessary} for a given task relative to a resource-matched classical alternative; distinguishing the two requires the sample-complexity and trainability arguments developed in Sections~\ref{sec:barren}, \ref{sec:symmetry}, and~\ref{sec:advantage}. Section~\ref{sec:notation} fixes notation. Section~\ref{sec:vqa} establishes the variational quantum algorithm formalism and gradient estimation used throughout. Section~\ref{sec:barren} develops the barren-plateau variance-scaling argument and surveys mitigation strategies, including restricted-connectivity, symmetry-restricted, and Hamiltonian-generated ansätze. Section~\ref{sec:qip} covers quantum image representations and their state-preparation complexity. Section~\ref{sec:qcnn} covers quantum convolutional neural networks. Section~\ref{sec:tensor} develops the tensor-network/MPS formalism in depth, including efficient loading conditions. Section~\ref{sec:qae} covers quantum autoencoders, dressed quantum circuits, and generative models. Section~\ref{sec:kernel} covers quantum kernel methods. Section~\ref{sec:symmetry} develops group representation theory and equivariance from first principles, worked through the dihedral group $D_4$. Section~\ref{sec:hamiltonian} covers physics-informed and Hamiltonian-based regularization, including the Heisenberg spin-chain symmetry hierarchy and its integrability. Section~\ref{sec:advantage} covers quantum advantage and dequantization. Section~\ref{sec:summary} gives a comparative summary table across the architecture families surveyed. Section~\ref{sec:challenges} discusses open challenges, and Section~\ref{sec:conclusion} concludes.

\section{Notation and Conventions}
\label{sec:notation}
Throughout, $n$ denotes the number of qubits, $\ket{0}=(1,0)^T$ and $\ket{1}=(0,1)^T$ the computational basis of a single qubit, and $X,Y,Z$ the Pauli operators. A quantum state on $n$ qubits is a unit vector $\ket{\psi}\in\mathcal{H}=(\mathbb{C}^2)^{\otimes n}$, or, more generally, a density operator $\rho\in\mathcal{B}(\mathcal{H})$ with $\rho\succeq 0$, $\Tr\rho=1$. A parameterized quantum circuit realizes a unitary $U(\bm{\theta})\in\mathrm{U}(2^n)$ for a parameter vector $\bm{\theta}\in\mathbb{R}^{m}$; we write $\rho(\bm{\theta})=U(\bm{\theta})\rho_0 U(\bm{\theta})^\dagger$ for the state prepared from a fixed reference $\rho_0$. Expectation values of an observable $O=O^\dagger$ are $\langle O\rangle_{\bm{\theta}} = \Tr[O\rho(\bm{\theta})]$. $G$ denotes a finite group with identity $e$, $|G|$ its order, and $\rho_{\mathcal{X}},\rho_{\mathcal{Y}}$ representations of $G$ on vector spaces $\mathcal{X},\mathcal{Y}$ (overloading $\rho$ for group representations and quantum density operators is standard in this literature; the correct reading is disambiguated by context in every instance below). $\chi\in\mathbb{N}$ denotes MPS bond dimension throughout Section~\ref{sec:tensor}, distinct from the susceptibility symbol $\mathcal{X}(t)$ used only in Section~\ref{sec:hamiltonian} discussion of training dynamics, which we denote with a calligraphic $\mathcal{X}$ specifically to avoid this collision.

\section{Variational Quantum Algorithms: Formalism and Gradient Estimation}
\label{sec:vqa}

\subsection{The variational framework}
A variational quantum algorithm (VQA) \cite{Cerezo2021} prepares a parameterized state
\begin{equation}
\ket{\psi(\bm{\theta})} = U(\bm{\theta})\ket{\psi_0} = \left(\prod_{l=1}^{L} W_l\, U_l(\theta_l)\right)\ket{0}^{\otimes n},
\end{equation}
where $\ket{\psi_0}=\ket{0}^{\otimes n}$ is a fixed reference state, each $U_l(\theta_l)=\exp(-i\theta_l G_l/2)$ is a single-parameter rotation generated by a Hermitian, typically involutory generator $G_l$ (so $G_l^2=\mathbb{1}$, e.g. a Pauli string), and each $W_l$ is a fixed entangling layer. A cost function
\begin{equation}
C(\bm{\theta}) = \bra{\psi(\bm{\theta})} H \ket{\psi(\bm{\theta})} = \Tr\left[H\,\rho(\bm{\theta})\right]
\end{equation}
is evaluated as the expectation value of a problem-defined observable $H$, and a classical optimizer updates $\bm{\theta}$ to minimize $C$. The variational quantum eigensolver (VQE) \cite{Peruzzo2014} is the canonical instance, with $H$ a molecular or spin-model Hamiltonian; a comprehensive treatment of VQE methodology and best practices is given in \cite{Tilly2022}, and its extension to quantum chemistry and materials science specifically in \cite{Bauer2020}. NISQ-era algorithm design more broadly is surveyed in \cite{Bharti2022}, with comprehensive treatments of the field's statistical-learning-theoretic foundations in \cite{Zaman2023,Wang2024QML}.

\subsection{The parameter-shift rule}
Because each $U_l(\theta_l)$ is generated by an involutory Hermitian operator with eigenvalues $\pm 1$, the expectation value $C(\bm{\theta})$ is exactly a single sinusoid in $\theta_l$ (holding all other parameters fixed): writing $G_l = P_+ - P_-$ for the projectors onto the $\pm1$ eigenspaces of $G_l$, one has $U_l(\theta_l) = \cos(\theta_l/2)\mathbb{1} - i\sin(\theta_l/2)G_l$, so $C(\theta_l)$ is a degree-one trigonometric polynomial in $\theta_l$ and is therefore exactly determined, and exactly differentiated, by evaluating it at two points:
\begin{equation}
\frac{\partial C}{\partial \theta_l} = \frac{1}{2}\left[C\!\left(\theta_l + \frac{\pi}{2}\right) - C\!\left(\theta_l - \frac{\pi}{2}\right)\right].
\label{eq:param_shift}
\end{equation}
Equation~\eqref{eq:param_shift} is the \emph{parameter-shift rule}: each partial derivative requires two additional circuit evaluations on the same hardware used for the forward pass, in contrast to a finite-difference estimator, which is biased at any finite step size and whose bias does not vanish under realistic hardware noise. This rule underlies essentially every trainable VQA and QML circuit discussed in this review.

\section{Trainability: Barren Plateaus}
\label{sec:barren}

\subsection{The variance-scaling argument}
McClean et al.\ \cite{McClean2018} established that for sufficiently expressive random parameterized circuits — formally, circuits whose layers form (at least) a 2-design over the unitary group $\mathrm{U}(2^n)$ — the variance of the cost-function gradient vanishes exponentially in the number of qubits $n$:
\begin{equation}
\mathrm{Var}\left[\frac{\partial C}{\partial\theta_l}\right] \in O\!\left(b^{-n}\right) \text{ for some } b>1, \qquad \mathbb{E}\left[\frac{\partial C}{\partial\theta_l}\right] = 0.
\label{eq:bp_variance}
\end{equation}
The proof strategy proceeds via Weingarten calculus: for $U$ Haar-random on $\mathrm{U}(2^n)$, moments $\mathbb{E}_U\left[U_{i_1 j_1}\cdots U_{i_k j_k}\overline{U_{i_1' j_1'}}\cdots\overline{U_{i_k' j_k'}}\right]$ of matrix entries reduce to a sum over pairings of the index sets weighted by Weingarten functions $\mathrm{Wg}(\sigma,2^n)$, which are rational functions of $2^n$ decaying as inverse powers of the Hilbert space dimension for fixed pairing structure $\sigma$. Applying this to $\mathbb{E}_U\left[\Tr(U\rho_0 U^\dagger [\partial_\theta U] \ldots)\right]$-type expressions that arise from differentiating $C(\bm{\theta})$ under the 2-design assumption yields both $\mathbb{E}[\partial_\theta C]=0$ (by a symmetry/cancellation argument under the Haar measure) and $\mathrm{Var}[\partial_\theta C]$ expressed as a ratio of Weingarten functions that is $O(4^{-n})$ for generic circuit bipartitions with global observables, tightening to $O(2^{-n})$ under some local-observable constructions — the precise base $b$ in Eq.~\eqref{eq:bp_variance} depends on the observable's locality and the circuit's light-cone structure, a refinement developed extensively in follow-up work \cite{Larocca2025}. Since $\mathrm{Var}[\partial_\theta C]\to 0$ exponentially while the number of measurement shots needed to resolve a quantity to fixed absolute precision scales as $\Theta(1/\mathrm{Var})$ by the standard shot-noise scaling of a sample mean, resolving a single gradient component above statistical noise requires exponentially many circuit executions — the \emph{barren plateau}. This is a statement about the \emph{landscape}, not the optimizer: no gradient-based, and per follow-up work many gradient-free, method can efficiently navigate a landscape satisfying Eq.~\eqref{eq:bp_variance} without exponentially many measurement shots.

\subsection{Sources of barren plateaus beyond circuit randomness}
Subsequent work established that circuit expressivity via approximate 2-design behavior is not the only route to Eq.~\eqref{eq:bp_variance}: global (rather than local, few-qubit-supported) observables induce vanishing gradients even for shallow, non-random circuits, since the variance of a global observable's expectation concentrates as a consequence of measure concentration on high-dimensional spheres (a manifestation of L\'evy's lemma) independent of circuit depth; hardware noise independently induces a distinct, depth-linear (rather than depth-independent-but-width-exponential) vanishing-gradient phenomenonsometimes termed a noise-induced barren plateau; and high circuit expressivity, even away from exact 2-design behavior, has been shown to correlate with vanishing gradient variance quite generally. This taxonomy — ansatz structure, observable locality, initial state, cost function choice, and hardware noise, each an independent axis along which a barren plateau can arise or be avoided — is systematically catalogued in the comprehensive review of Larocca et al.\ \cite{Larocca2025}, with concrete mitigation strategies (initialization schemes, layerwise training, identity-block initialization, and the structural constraints discussed next) surveyed in \cite{Cunningham2024}.

\subsection{Architectural mitigation}
Three classes of architectural mitigation are directly relevant to vision applications. First, \emph{restricted-connectivity} ansätze: the quantum convolutional neural network (Section~\ref{sec:qcnn}) is the leading example, in which Pesah et al.\ \cite{Pesah2021} prove that the alternating convolution/pooling structure, which discards half the qubits at each pooling layer, bounds the effective circuit depth seen by any single output qubit logarithmically in $n$, giving polynomially rather than exponentially vanishing gradients. Second, \emph{symmetry-restricted} ansätze (Section~\ref{sec:symmetry}), which restrict the parameterized unitary to a subgroup or a dynamical Lie subalgebra small enough to avoid generic 2-design behavior \cite{West2024,Nguyen2024EQNN,Schatzki2024}. Third, \emph{Hamiltonian-generated} ansätze, in which each layer's generator $G_l$ is drawn from a fixed physical Hamiltonian (e.g. Heisenberg exchange terms, Section~\ref{sec:hamiltonian}) rather than an unconstrained Pauli string — Park and Killoran \cite{Park2024} prove that Hamiltonian variational ansätze of this form retain polynomially bounded gradient variance under mild conditions on the generating Hamiltonian's locality structure, a result of direct relevance to any architecture that uses a physically structured generator set rather than a generic hardware-efficient ansatz. A fourth, more radical route uses mid-circuit measurement and classical feedforward to construct circuits that remain expressive while being provably free of the barren-plateau obstruction \cite{Deshpande2024}.

\section{Quantum Representations of Images}
\label{sec:qip}

\subsection{Amplitude encodings: FRQI and NEQR}
The Flexible Representation of Quantum Images (FRQI) \cite{Le2011} encodes a $2^n\times 2^n$ grayscale image as
\begin{equation}
\ket{I} = \frac{1}{2^n}\sum_{i=0}^{2^{2n}-1}\left(\cos\theta_i\ket{0} + \sin\theta_i\ket{1}\right)\otimes\ket{i},
\end{equation}
using $2n+1$ qubits ($2n$ for the position register $\ket{i}$, one color qubit), where $\theta_i\in[0,\pi/2]$ encodes normalized grayscale intensity as a rotation amplitude entangled with position. Because intensity is encoded in an \emph{amplitude}, recovering any single pixel's value to precision $\varepsilon$ requires $\Theta(1/\varepsilon^2)$ repeated measurements by the standard sample-mean shot-noise scaling, and the general FRQI state-preparation circuit requires $O(4^n)$ elementary gates in the worst case (one controlled rotation per position basis state, uncompressed). The Novel Enhanced Quantum Representation (NEQR) \cite{Zhang2013} instead encodes intensity in a \emph{computational basis state} of a bank of $q$ qubits (typically $q=8$ for 256 gray levels) using $2n+q$ total qubits,
\begin{equation}
\ket{I} = \frac{1}{2^n}\sum_{i=0}^{2^{2n}-1}\ket{c_i}\otimes\ket{i}, \qquad \ket{c_i}=\ket{c^{q-1}_i\cdots c^{0}_i}, \quad c_i = \sum_{k=0}^{q-1} c_i^k 2^k,
\end{equation}
which permits exact intensity recovery from a single noiseless measurement and admits more efficient pixel-level operations such as translation and edge detection \cite{Zhou2017,Yao2017}, at the cost of $O(q\cdot 2^{2n})$ two-qubit gates for generic (incompressible) images in the naive construction. Both encodings scale the position register with $\log_2$ of the pixel count, which for realistic image resolutions (tens of thousands of pixels) already exceeds current and near-term qubit counts; this is the primary reason essentially every hybrid architecture surveyed here restricts direct quantum encoding to small patches (typically $\lesssim 64$ pixels) and relegates full-image structure to classical pre- or post-processing. Comparative treatments of FRQI-versus-NEQR trade-offs, additional pixel-representation variants, and DCT-domain quantum encodings that reduce resource cost on natural (as opposed to synthetic worst-case) images are given in \cite{QIPSystematicReview2025,QuantumImageRepClassification2025}, with \cite{Chakraborty2018} surveying the broader set of open algorithmic and hardware challenges specific to quantum image processing as a subfield distinct from quantum machine learning generally.

\subsection{Datasets and evaluation practice}
Because the qubit and gate-count constraints of Section~\ref{sec:qip}.1 restrict essentially every architecture surveyed in this review to small or aggressively downsampled inputs, the field has converged on a narrow, recurring set of benchmark datasets rather than the high-resolution, large-scale benchmarks standard in classical computer vision. MNIST and Fashion-MNIST ($28\times28$, or further downsampled/PCA-reduced to match qubit count) dominate the QCNN and dressed-circuit literature (Section~\ref{sec:qcnn}, \cite{Hur2022,Cong2019,Mari2020}); CIFAR-10 at native or reduced resolution appears across kernel, tensor-network, and hybrid-autoencoder studies (Section~\ref{sec:kernel}, \ref{sec:tensor}, \cite{Slabbert2024,Guala2023}); and medical-imaging benchmarks (chest X-ray, histopathology, and similar modalities) appear disproportionately often in the hybrid-pipeline literature relative to their prevalence in classical vision benchmarking, plausibly because the classification tasks on such datasets are frequently lower-dimensional and higher-contrast than natural-image benchmarks, and therefore more forgiving of the aggressive downsampling every quantum encoding in Section~\ref{sec:qip} requires. This benchmark concentration matters for interpreting the empirical results throughout this review: a reported accuracy competitive with a classical baseline on $28\times28$ MNIST does not, without further argument, extrapolate to a claim about performance at the resolutions used in mainstream classical vision benchmarking, and the resource-matching discipline discussed in Section~\ref{sec:challenges} needs to hold not just for parameter count but for the effective information content of the (typically heavily compressed) input each side of a comparison actually receives.

\section{Quantum Convolutional Neural Networks}
\label{sec:qcnn}

\subsection{The MERA-derived architecture}
The quantum convolutional neural network (QCNN) \cite{Cong2019} is a circuit ansatz directly inspired by the multiscale entanglement renormalization ansatz (MERA) from condensed-matter physics, alternating two structurally distinct layer types over $\log_2 n$ stages applied to $n$ input qubits: a \emph{convolutional} layer applies a single, translationally-shared two-qubit (or few-qubit) unitary $U_{\text{conv}}$ across all nearest-neighbor pairs, and a \emph{pooling} layer measures a subset (typically half) of the remaining qubits and applies a single-qubit unitary $V(\bm{m})$ to each surviving neighbor, conditioned on the classical measurement outcome $\bm{m}$ of its pooled partner. After $\log_2 n$ stages, a constant number of qubits remain, and a final fully-connected layer of gates on the surviving qubits produces the output via measurement of a designated observable. Because a fixed, shared set of parameters $(U_{\text{conv}},V)$ is applied at every stage, the total parameter count scales as $O(\log n)$ rather than $O(n)$ or $O(n^2)$, and because pooling geometrically halves the qubit count at each stage, the circuit depth relevant to any single output qubit's light cone is $O(\log n)$ — precisely the structural property Pesah et al.\ \cite{Pesah2021} exploit to prove QCNNs are free of the exponential barren-plateau scaling of Eq.~\eqref{eq:bp_variance}, retaining only a polynomial decay in $n$.

\subsection{Empirical benchmarking}
Hur, Kim, and Park \cite{Hur2022} give the most extensive published empirical comparison of QCNN design choices for classical image data specifically, systematically varying ansatz structure, quantum data-encoding method (amplitude, angle, and basis encodings), classical pre-processing (raw pixels versus PCA-reduced features), cost function, and classical optimizer across binary and multi-class subsets of MNIST and Fashion-MNIST. Their central empirical finding is that the choice of \emph{encoding} strategy and classical pre-processing dominates final accuracy more than the choice of ansatz or pooling variant at matched parameter count, suggesting — consistent with the power-of-data framing of Section~\ref{sec:advantage} — that how classical information enters the quantum circuit is at least as consequential as the circuit's internal entangling structure. Complementary reviews and tutorials on the QCNN circuit construction, deterministic-versus-measurement-based pooling variants, and integration into hierarchical feature-extraction pipelines further systematize this design space; digital-analog and interaction-layer variants of the base QCNN architecture have subsequently been proposed to further reduce circuit depth on specific hardware connectivity graphs.

\section{Tensor Networks for Machine Learning}
\label{sec:tensor}

\subsection{The matrix product state and singular value decomposition}
A generic pure state of $N$ qubits, $\ket{\psi}=\sum_{s_1,\ldots,s_N} c_{s_1\cdots s_N}\ket{s_1}\cdots\ket{s_N}$, has a coefficient tensor $c$ with $2^N$ entries. The matrix product state (MPS) ansatz factorizes this tensor as a chain of smaller tensors,
\begin{equation}
c_{s_1 s_2\cdots s_N} = \sum_{\alpha_1,\ldots,\alpha_{N-1}} A^{s_1}_{\alpha_1} A^{s_2}_{\alpha_1\alpha_2}\, A^{s_3}_{\alpha_2\alpha_3}\cdots A^{s_N}_{\alpha_{N-1}},
\label{eq:mps}
\end{equation}
where each $A^{s_k}$ is a matrix of dimension $\chi_{k-1}\times\chi_k$, the bond dimensions $\chi_k$ bound the entanglement representable across the cut between sites $k$ and $k+1$, and any state can be brought to this form exactly (with $\chi_k \le 2^{\min(k,N-k)}$) via a sequence of $N-1$ singular value decompositions sweeping across the chain: reshaping $c$ into a matrix $C_{(s_1\cdots s_k),(s_{k+1}\cdots s_N)}$ across the cut at site $k$, the SVD $C = U\Sigma V^\dagger$ yields a canonical bipartition with Schmidt coefficients on the diagonal of $\Sigma$. Truncating each SVD to retain only the $\chi \ll 2^{\min(k,N-k)}$ largest singular values yields a controlled, variational approximation whose squared truncation error is exactly the discarded Schmidt-weight, $\sum_{k>\chi}\lambda_k^2$ (for singular values normalized so $\sum_k \lambda_k^2=1$); this truncation is precisely the mechanism by which patch-based architectures compress a $2^{12}$-dimensional 12-site patch state into a bond-dimension-4 MPS. The bipartite entanglement (von Neumann) entropy across a cut is computed directly from the normalized singular values $p_k = \lambda_k^2$,
\begin{equation}
S = -\sum_{k=1}^{\chi} p_k \log p_k \ \le\ \log_2\chi,
\label{eq:ee}
\end{equation}
which Latorre \cite{Latorre2005} first identified as the quantity linking image compressibility to entanglement: an MPS representation of a (suitably encoded) classical image compresses well under bond-dimension truncation exactly when its entanglement entropy profile — equivalently, its Schmidt spectrum decay — is concentrated on a small effective number of terms, a criterion satisfied by locally correlated, low-frequency image content and violated by highly non-local, texture-rich content. Liu et al.\ \cite{Liu2021} operationalize the same quantity for feature selection, ranking individual pixels by their contribution to the site-resolved entanglement entropy and showing that the great majority of pixels in typical natural images can be discarded without loss of downstream classification accuracy.

\subsection{Efficient loading: the Fourier-spectrum condition}
A separate and practically crucial question is not whether a compressed MPS representation exists, but whether it can be \emph{loaded onto a quantum circuit} with a gate count that scales favorably. Jobst et al.\ \cite{Jobst2024} address this directly for classical image data: they show that images whose two-dimensional discrete Fourier spectrum decays sufficiently rapidly with spatial frequency admit an MPS representation with Schmidt rank (bond dimension) growing only polylogarithmically, rather than exponentially, in the number of encoded pixels, and give an explicit, analytically bounded procedure for constructing the corresponding quantum circuit from the Fourier coefficients directly, rather than via a generic (and generally exponential-depth) arbitrary-state-preparation routine. Their empirical validation on natural-image benchmarks (Imagenette, DIV2K) further finds that simple, sequential (as opposed to more expressive, densely entangled) circuit ansätze suffice to match the reconstruction fidelity of richer ansätze once the input has been compressed via this Fourier-decay-informed MPS truncation, an observation consistent with the broader pattern in Section~\ref{sec:advantage} that architectural expressivity beyond what the data's own structure requires does not translate into improved empirical performance.

\subsection{Tensor networks as a supervised and generative learning primitive}
Independently of any quantum hardware execution, the MPS ansatz of Eq.~\eqref{eq:mps} is itself a trainable machine learning model: Stoudenmire and Schwab \cite{Stoudenmire2016} parameterize a linear classifier as an MPS acting on a per-pixel local feature map (e.g. $\phi(x_i)=(\cos(\pi x_i/2),\sin(\pi x_i/2))$ tensored across pixels), trained via a DMRG-style two-site sweep in which adjacent tensors $A^{s_k}A^{s_{k+1}}$ are merged into a single two-site tensor, updated by gradient descent (or an eigensolver, in the original DMRG context) against the training loss, and then re-split via SVD truncation back to bond dimension $\chi$ before sweeping to the next bond — and report below 1\% test error on MNIST digit classification, establishing that an appreciable fraction of MPS-based vision performance is attributable to the tensor-network inductive bias itself, independent of any quantum circuit that might later measure observables from a related state. Han et al.\ \cite{Han2018} extend the same ansatz to unsupervised generative modeling, exploiting the MPS's built-in exact-sampling procedure (the canonical form used for the SVD truncation above also supports sequential, exact conditional sampling site-by-site) to draw new samples without the approximate MCMC procedures generative energy-based models typically require. Guala et al.\ \cite{Guala2023} translate MPS/tree-tensor-network classifiers into variational quantum circuits proper, benchmarking image classification with circuit-cutting techniques to scale beyond the qubit counts available on a single device, and Wang et al.\ \cite{Wang2023TN} survey the broader convergence of tensor-network structure with mainstream classical architectures (CNNs, RNNs, transformers, graph networks) under the umbrella of ``tensorial neural networks.'' Cirac et al.\ \cite{Cirac2021} give the definitive review of MPS and its two-dimensional generalization, the projected entangled pair state (PEPS), as representations of quantum many-body states, the formal setting the machine learning constructions above are borrowed from. Tensor-network quantum circuits have also been applied outside vision proper, e.g. to particle-physics event classification \cite{Araz2022}, indicating the inductive bias generalizes across structured, locally-correlated data more broadly than images specifically.

\section{Quantum Autoencoders, Dressed Circuits, and Generative Models}
\label{sec:qae}

\subsection{Quantum autoencoders}
The original quantum autoencoder (QAE) \cite{Romero2017} compresses an ensemble of $n$-qubit states $\{\ket{\psi_i}\}$ into $n-k$ ``trash'' qubits plus $k$ ``latent'' qubits by training a parameterized unitary $U(\bm{\theta})$ to maximize the fidelity between the trash-register reduced state and a fixed reference (conventionally $\ket{0}^{\otimes(n-k)}$), typically estimated via a swap test:
\begin{equation}
\bm{\theta}^\star = \arg\max_{\bm{\theta}}\ \frac{1}{|\mathcal{D}|}\sum_{i\in\mathcal{D}} \bra{0}^{\otimes(n-k)}\Tr_{\text{latent}}\!\left[U(\bm{\theta})\ket{\psi_i}\!\bra{\psi_i}U(\bm{\theta})^\dagger\right]\ket{0}^{\otimes(n-k)}.
\label{eq:qae_fidelity}
\end{equation}
This objective compresses genuinely \emph{quantum} data; when applied to classically encoded images, Eq.~\eqref{eq:qae_fidelity} becomes a quantum-native analogue of a classical autoencoder's reconstruction loss. Hybrid variants attach a classical decoder directly to measured amplitudes or observables rather than training the trash-qubit fidelity objective at all \cite{WangQAECompression2024,QAEClassification2025}, and a related but architecturally distinct proposal compresses information encoded in the circuit itself (its gate sequence) rather than an input quantum state \cite{QCAE2023}.

\subsection{Dressed quantum circuits and hybrid transfer learning}
A structurally distinct and empirically influential pattern, introduced by Mari et al.\ \cite{Mari2020}, embeds a small variational quantum circuit \emph{between} two classical networks rather than at either end of the pipeline: a classical pre-processing network $f_{\text{pre}}$ (frequently a pre-trained, frozen backbone such as a ResNet) reduces a high-dimensional input to a $k$-dimensional feature vector matching the available qubit count, which is angle- or amplitude-encoded and processed by a variational circuit $U(\bm{\theta})$, whose measured expectation values are handed to a classical post-processing network $f_{\text{post}}$ producing the final output,
\begin{equation}
y = f_{\text{post}}\Bigl(\bigl\langle O_1\bigr\rangle_{\bm{\theta}},\ldots,\bigl\langle O_k\bigr\rangle_{\bm{\theta}}\Bigr), \qquad \langle O_j\rangle_{\bm{\theta}} = \bra{0}^{\otimes k} U_{\text{enc}}(f_{\text{pre}}(x))^\dagger U(\bm{\theta})^\dagger O_j\, U(\bm{\theta}) U_{\text{enc}}(f_{\text{pre}}(x))\ket{0}^{\otimes k},
\label{eq:dressed}
\end{equation}
with $f_{\text{pre}}$, $\bm{\theta}$, and $f_{\text{post}}$ trained jointly end-to-end via the parameter-shift rule (Eq.~\eqref{eq:param_shift}) for the quantum parameters and ordinary backpropagation for the classical ones. This ``dressed quantum circuit'' pattern is architecturally the pattern that essentially all vision-scale hybrid pipelines surveyed in this review instantiate in one form or another (Sections~\ref{sec:qcnn}--\ref{sec:tensor} being partial exceptions in that they apply the quantum structure closer to the raw input); Eq.~\eqref{eq:dressed}'s central design freedom — how much representational work is delegated to $f_{\text{pre}}/f_{\text{post}}$ versus to $U(\bm{\theta})$ — is exactly the axis the power-of-data framework of Section~\ref{sec:advantage} formalizes as the determinant of whether the quantum component contributes anything a resource-matched classical alternative could not.

\subsection{Quantum and hybrid generative models}
Quantum generative adversarial networks (qGANs) \cite{Zoufal2019} replace one or both of the generator and discriminator of a classical GAN with a parameterized quantum circuit, trained via the standard minimax objective
\begin{equation}
\min_{G} \max_{D}\ \mathbb{E}_{x\sim p_{\text{data}}}\left[\log D(x)\right] + \mathbb{E}_{z\sim p_z}\left[\log\bigl(1-D(G(z))\bigr)\right],
\end{equation}
with $G$ (and optionally $D$) implemented as PQCs whose output is read out via measured expectation values or sampled bitstrings; Huang et al.\ \cite{HuangExpGAN2020} demonstrate this experimentally for handwritten-digit generation on superconducting hardware, and subsequent hybrid variants \cite{HybridGANHighRes2022,MosaiQ2023} target higher resolution and NISQ-noise robustness respectively, the latter reporting parameter counts several orders of magnitude smaller than a comparable classical GAN for matched output fidelity — a genuine expressivity-per-parameter claim distinct from, and not automatically implying, a sample- or compute-complexity advantage (Section~\ref{sec:advantage}). Quantum denoising diffusion models \cite{QuantumDiffusion2024} and hybrid latent-diffusion variants for medical imaging \cite{HybridLatentDiffusion2025} extend the same hybridization pattern to the diffusion generative framework.

\section{Quantum Kernel Methods}
\label{sec:kernel}
A quantum feature map $\phi:\mathcal{X}\to\mathcal{H}$ embeds a classical input $x$ into a quantum state $\ket{\phi(x)} = U_\phi(x)\ket{0}^{\otimes n}$ prepared by a data-dependent circuit $U_\phi(x)$; the induced quantum kernel is the squared state overlap,
\begin{equation}
\kappa(x,x') = \left|\braket{\phi(x)}{\phi(x')}\right|^2 = \left|\bra{0}^{\otimes n} U_\phi(x)^\dagger U_\phi(x') \ket{0}^{\otimes n}\right|^2,
\label{eq:qkernel}
\end{equation}
which is estimated on hardware via the same overlap circuit and a finite number of measurement shots (with shot-noise-limited estimation error scaling as in Section~\ref{sec:qip}), and is then handed to a classical kernel method exactly as any classical kernel would be. Havlíček et al.\ \cite{Havlicek2019} introduce this construction and show, on a synthetic dataset engineered to be classically hard, that a quantum feature map based on a circuit whose classical simulation is conjectured intractable yields a separable kernel where classical kernels do not; applications to real image data include quantum kernels for manufacturing-defect detection on superconducting hardware \cite{QuantumKernelDefect2022} and hybrid classical-autoencoder-plus-quantum-SVM pipelines for image classification \cite{Slabbert2024}. The reproducing-kernel-Hilbert-space (RKHS) view of Eq.~\eqref{eq:qkernel} is what connects this section directly to the dequantization literature of Section~\ref{sec:advantage}: because $\kappa$ is, after estimation, an ordinary positive-semidefinite kernel matrix, any advantage of a quantum feature map over a classical one must be located in the \emph{kernel values themselves} being hard to compute or approximate classically, not merely in the downstream classical learning algorithm that consumes them.

\section{Symmetry and Equivariance}
\label{sec:symmetry}

\subsection{Group actions and equivariant maps}
Let $G$ be a finite group acting on an input space $\mathcal{X}$ via a representation $\rho_{\mathcal{X}}:G\to\mathrm{GL}(\mathcal{X})$ and on an output space $\mathcal{Y}$ via $\rho_{\mathcal{Y}}:G\to\mathrm{GL}(\mathcal{Y})$. A map $f:\mathcal{X}\to\mathcal{Y}$ is \emph{$G$-equivariant} if
\begin{equation}
f\bigl(\rho_{\mathcal{X}}(g)\,x\bigr) = \rho_{\mathcal{Y}}(g)\, f(x) \qquad \text{for all } g\in G,\ x\in\mathcal{X},
\label{eq:equivariance}
\end{equation}
and \emph{$G$-invariant} in the special case $\rho_{\mathcal{Y}}(g)=\mathrm{id}$ for all $g$. Group-equivariant convolutional networks (G-CNNs) \cite{Cohen2016} generalize ordinary translation-equivariant convolution to Eq.~\eqref{eq:equivariance} for arbitrary finite (and, in later work, continuous or steerable \cite{Weiler2019}) groups $G$, and this line of work is now organized, together with graph- and manifold-structured generalizations, under geometric deep learning \cite{Bronstein2021}. The architectural payoff of Eq.~\eqref{eq:equivariance} is a strict reduction in the effective hypothesis class relative to an unconstrained model of the same raw parameter count, which reduces sample complexity and, as discussed in Section~\ref{sec:barren}, is one of the few systematically understood routes to bounding rather than merely observing a variational quantum circuit's trainability.

\subsection{Worked example: the dihedral group $D_4$}
The dihedral group of the square, $D_4 = \langle r,s \mid r^4=s^2=(sr)^2=e\rangle$, has order $8$: four rotations $\{e,r,r^2,r^3\}$ (by $0^\circ,90^\circ,180^\circ,270^\circ$) and four reflections $\{s,rs,r^2s,r^3s\}$. $D_4$ has five conjugacy classes and five irreducible representations — four one-dimensional ($A_1$ trivial, $A_2$, $B_1$, $B_2$) and one two-dimensional ($E$) — with character table
\begin{center}
\begin{tabular}{lccccc}
\toprule
& $e$ & $r,r^3$ & $r^2$ & $s,r^2s$ & $rs,r^3s$ \\
\midrule
$A_1$ & 1 & 1 & 1 & 1 & 1 \\
$A_2$ & 1 & 1 & 1 & $-1$ & $-1$ \\
$B_1$ & 1 & $-1$ & 1 & 1 & $-1$ \\
$B_2$ & 1 & $-1$ & 1 & $-1$ & 1 \\
$E$ & 2 & 0 & $-2$ & 0 & 0 \\
\bottomrule
\end{tabular}
\end{center}
so that any linear representation $V$ of $D_4$ decomposes as $V \cong m_{A_1} A_1 \oplus m_{A_2} A_2 \oplus m_{B_1} B_1 \oplus m_{B_2} B_2 \oplus m_E E$, with multiplicities given by the character-orthogonality inner product $m_i = \frac{1}{|D_4|}\sum_{g\in D_4}\overline{\chi_i(g)}\,\chi_V(g)$. For a feature map $\Phi:\mathcal{X}\to\mathbb{R}^d$ on a square patch grid that is \emph{not} already equivariant, the canonical construction of an exactly $D_4$-equivariant map from it is Reynolds/group averaging — a projector, in the sense of Eq.~\eqref{eq:d4_decomp_ref}, onto the invariant-plus-equivariant subspace:
\begin{equation}
\Phi_{D_4}(x) = \frac{1}{|D_4|}\sum_{g\in D_4} \rho_{\mathcal{Y}}(g)^{-1}\, \Phi\bigl(\rho_{\mathcal{X}}(g)\, x\bigr) = \frac{1}{8}\sum_{g\in D_4} P_g^{-1}\Phi(gx),
\label{eq:reynolds}
\end{equation}
which satisfies Eq.~\eqref{eq:equivariance} \emph{exactly}, up to floating-point arithmetic, for \emph{any} base map $\Phi$: writing $\Phi = \sum_i \Phi_i$ for its decomposition into isotypic components under the action of Eq.~\eqref{eq:d4_decomp_ref}, the group-average operator in Eq.~\eqref{eq:reynolds} is precisely the orthogonal projector onto the trivial isotypic component in the argument and identity on the value, by the standard orthogonality-of-characters argument used to derive the multiplicity formula above; every non-invariant isotypic component is annihilated by the sum. This is worth emphasizing as a mathematical point with direct empirical consequence: a feature map merely \emph{motivated} by a symmetry (e.g. built from a positional encoding that respects a square grid's geometry) is, in general, \emph{not} equivariant in the sense of Eq.~\eqref{eq:equivariance} unless it is explicitly symmetrized via a construction such as Eq.~\eqref{eq:reynolds} or is provably invariant by the structure of its generators alone; geometric motivation and enforced equivariance are mathematically distinct properties, and only the latter is a guarantee.
\begin{equation}
V \cong m_{A_1} A_1 \oplus m_{A_2} A_2 \oplus m_{B_1} B_1 \oplus m_{B_2} B_2 \oplus m_E E
\label{eq:d4_decomp_ref}
\end{equation}

\subsection{Equivariant variational quantum circuits}
Meyer et al.\ \cite{Meyer2023} construct the quantum analogue of Section~\ref{sec:symmetry}.2 by drawing PQC generators $G_l$ from the commutant of a unitary representation of $G$ on the Hilbert space, so that $U(\bm{\theta})\rho_{\mathcal{H}}(g) = \rho_{\mathcal{H}}(g)U(\bm{\theta})$ for all $g\in G$ by construction rather than by training; Nguyen et al.\ \cite{Nguyen2024EQNN} generalize this into a full representation-theoretic framework for equivariant quantum neural networks (EQNNs) applicable to arbitrary symmetry groups, demonstrating $\mathrm{SU}(2)$-equivariant quantum convolutional networks on the bond-alternating Heisenberg model — directly connecting the symmetry formalism of this section to the Heisenberg-model formalism of Section~\ref{sec:hamiltonian}. Schatzki et al.\ \cite{Schatzki2024} prove trainability and generalization guarantees specifically for permutation-equivariant QNNs, and Sauvage et al.\ \cite{Sauvage2024} construct circuits equivariant under the automorphism group of a problem Hamiltonian directly, showing faster training as a direct consequence. For vision tasks specifically, West et al.\ construct rotation-equivariant \cite{West2024} and reflection-equivariant \cite{West2023} variational quantum classifiers, reporting both provable freedom from barren plateaus and improved empirical performance on image classification relative to generic ansätze of matched size, with a subsequent rotation-\emph{invariant} construction \cite{West2025} extending the same principle further; a published Comment on the original rotation-equivariant construction flags a technical inconsistency in its dynamical-Lie-algebra argument, illustrating that even within this comparatively mathematically mature subfield, claims of provable trainability warrant independent verification rather than citation at face value. Skolik et al.\ \cite{Skolik2023} apply the same equivariant-circuit-design principle to graph-structured data via node-permutation equivariance, indicating the construction generalizes across data topologies, not solely image patch grids.

\section{Physics-Informed and Hamiltonian-Based Regularization}
\label{sec:hamiltonian}

\subsection{The anisotropic Heisenberg model and its symmetry hierarchy}
The anisotropic Heisenberg (XYZ) Hamiltonian on a chain of $n_{\text{bonds}}$ nearest-neighbor bonds,
\begin{equation}
H_{\text{Heis}} = \sum_{i=0}^{n_{\text{bonds}}-1} \left(J_x^{(i)} X_i X_{i+1} + J_y^{(i)} Y_i Y_{i+1} + J_z^{(i)} Z_i Z_{i+1}\right),
\label{eq:heisenberg}
\end{equation}
exhibits a symmetry hierarchy governed by the relations among $J_x^{(i)},J_y^{(i)},J_z^{(i)}$ \cite{Craps2024,Franchini2017}: for $J_x=J_y=J_z$ (isotropic, XXX), Eq.~\eqref{eq:heisenberg} commutes with the full $\mathrm{SU}(2)$ spin-rotation group generated by total spin $\bm{S}=\sum_i \bm{\sigma}_i/2$; for $J_x=J_y\neq J_z$ (XXZ), this reduces to the $\mathrm{U}(1)$ subgroup generated by total $S^z=\sum_i Z_i/2$; and for generic $J_x\neq J_y\neq J_z$ (XYZ), continuous symmetry is broken entirely, leaving only discrete $\mathbb{Z}_2\times\mathbb{Z}_2$ symmetry generated by $\pi$-rotations about two orthogonal spin axes.

\subsection{Integrability and the algebraic Bethe ansatz}
The XXZ chain is additionally exactly solvable via the algebraic Bethe ansatz \cite{Franchini2017}, a structure worth sketching because it is what makes the Heisenberg model, among possible physically motivated Hamiltonians, an unusually well-characterized choice of inductive bias. The construction proceeds via a Lax operator $L_{0,i}(\lambda)$ acting on an auxiliary space $0$ and physical site $i$, satisfying the Yang--Baxter equation
\begin{equation}
R_{12}(\lambda-\mu)\,L_{1}(\lambda)\,L_{2}(\mu) = L_{2}(\mu)\,L_{1}(\lambda)\,R_{12}(\lambda-\mu),
\end{equation}
for an $R$-matrix $R_{12}(\lambda)$ satisfying the Yang--Baxter equation in its own right; the monodromy matrix $T_0(\lambda)=L_{0,N}(\lambda)\cdots L_{0,1}(\lambda)$ built from the Lax operators across all $N$ sites yields, via its trace over the auxiliary space, a family of mutually commuting transfer matrices $\tau(\lambda)=\Tr_0 T_0(\lambda)$ whose expansion in $\lambda$ generates the Heisenberg Hamiltonian itself together with an extensive tower of higher conserved charges — the hallmark of quantum integrability. Eigenstates of $\tau(\lambda)$, and hence of $H_{\text{Heis}}$, are constructed algebraically by acting with off-diagonal monodromy-matrix elements on a reference (ferromagnetic) state, with the resulting spectrum and eigenstate structure governed by a closed system of algebraic (Bethe) equations for a set of rapidity parameters, rather than requiring numerical diagonalization. This is the sense in which the coupling structure of Eq.~\eqref{eq:heisenberg} is not merely a plausible physical prior but a mathematically exhaustively characterized one: when a learnable-coupling regularizer of this form is fit to data, the couplings' drift toward the isotropic, XXZ, or generic regime is a movement along a symmetry hierarchy whose representation theory, spectrum, and correlation functions are independently and exactly known, providing an interpretability anchor a generically parameterized energy function does not have.

\subsection{Parameterized Hamiltonian learning and physics-informed circuits}
A distinct line of work treats the Hamiltonian itself, rather than a downstream task loss, as the primary learned object: Shi et al.\ \cite{Shi2022} formulate parameterized Hamiltonian learning (PHL) as recovering an unknown Hamiltonian's structure directly from measurement data via a trainable quantum circuit whose parameters are the Hamiltonian's own coupling constants, applied there to image-segmentation-style tasks via an energy-landscape decision boundary rather than a reconstruction objective; physics-informed neural network (PINN) approaches to the same inverse problem \cite{Chinni2023} instead embed the Schr\"odinger or Heisenberg equation of motion directly into a classical network's loss function and recover two-qubit Hamiltonian parameters from real IBM hardware measurement data via gradient descent on the resulting physics-constrained objective. Both differ structurally from a Hamiltonian used as an \emph{auxiliary regularizer} alongside an independent primary loss (the pattern of Section~\ref{sec:vqa}): in PHL and Hamiltonian-recovery PINNs, the Hamiltonian's parameters are the entire object of inference and correctly recovering them is the success criterion, whereas an auxiliary regularizer's couplings are free to settle wherever a downstream task's gradient signal pushes them, without any guarantee that the resulting couplings correspond to a physically meaningful energy landscape at all — an empirical question, not a structural one, and precisely the kind of question a resource-matched ablation (rather than the architectural motivation alone) is needed to resolve.

\section{Quantum Advantage, Dequantization, and the Power of Data}
\label{sec:advantage}

\subsection{Dequantization}
A dequantization result exhibits a classical algorithm that matches a quantum algorithm's asymptotic performance on a specific task given comparable input access, thereby removing that task from the set of demonstrated quantum speedups. Tang's dequantization of the Kerenidis--Prakash quantum recommendation algorithm \cite{Tang2019} is the canonical example: given the same $\ell^2$-norm sampling access to a matrix that the quantum algorithm assumes, a classical randomized algorithm reproduces the quantum algorithm's polylogarithmic dependence on matrix dimension, at the cost of polynomial (rather than polylogarithmic) dependence on other problem parameters such as matrix rank and condition number — an exponential separation in one parameter can therefore coexist with no separation once the full parameter dependence is accounted for. Cotler et al.\ \cite{Cotler2021} sharpen the picture by distinguishing input-access models precisely: a classical algorithm with the same sample-and-query access a quantum algorithm has can, for a range of learning tasks, match or beat it, whereas genuine separations survive only under more restrictive classical access models. Gyurik and Dunjko \cite{Gyurik2023} give the complementary positive result, constructing explicit learning problems with provable, non-dequantizable exponential separations under cryptographic hardness assumptions — establishing that dequantization is not universal, but that its scope must be checked task-by-task rather than assumed away.

\subsection{The power of data}
Huang et al.\ \cite{Huang2021PowerData} reframe the advantage question in explicitly data-dependent terms: whether a quantum model outperforms a classical one on a given task is a property of the joint distribution over inputs, labels, and the specific feature map used, not an intrinsic property of ``quantum'' versus ``classical'' model classes in the abstract. They construct a projected quantum kernel that provably improves generalization when training data carries structure a classical model cannot efficiently access, and give a practical geometric test — computable from a modest amount of data before any large-scale training run — for whether a given task falls into the regime where a quantum model can be expected to outperform the best classical alternative. Applied to the constructions surveyed above (quantum kernels, Section~\ref{sec:kernel}; symmetry-restricted circuits, Section~\ref{sec:symmetry}; MPS/tensor-network feature maps, Section~\ref{sec:tensor}; dressed circuits, Section~\ref{sec:qae}), this framing implies architectural motivation is necessary but not sufficient evidence for an advantage: the target task's structure must independently require what the architecture provides, precisely the criterion formalized by the classical-shadows / few-measurement-estimation literature \cite{Huang2020Shadows} for how much classical post-processing can be extracted from a fixed measurement budget before any advantage claim about the measurement strategy itself is warranted.

\section{Comparative Summary}
\label{sec:summary}
Table~\ref{tab:summary} collects the architecture families and representative works surveyed above, organized by their primary mechanism, typical qubit-count regime, and the task-level claim each makes, to make explicit that ``quantum vision'' is not one architecture but at least seven structurally distinct research programs that happen to share a common hardware substrate.

\begin{table}[h]
\centering
\caption{Architecture families surveyed, representative works, and their primary mechanism.}
\label{tab:summary}
\scriptsize
\begin{tabular}{p{0.19\textwidth}p{0.20\textwidth}p{0.10\textwidth}p{0.42\textwidth}}
\toprule
Family & Representative works & Typical $n$ & Primary mechanism \\
\midrule
Quantum image encodings & \cite{Le2011,Zhang2013} & $2n{+}q$ & Exact/amplitude pixel encoding into basis states or rotation angles. \\
QCNN & \cite{Cong2019,Pesah2021,Hur2022} & $\le 16$ & MERA-derived alternating conv./pooling; $O(\log n)$ depth, parameter count. \\
Tensor-network feature maps & \cite{Stoudenmire2016,Han2018,Guala2023,Jobst2024} & $\le 16$ & MPS/SVD-truncated compression; Fourier-decay-conditioned efficient loading. \\
Quantum autoencoders & \cite{Romero2017,QCAE2023} & varies & Swap-test trash-qubit fidelity training objective. \\
Dressed circuits / transfer learning & \cite{Mari2020} & $\le 8$ & Classical pre-/post-network with a thin quantum layer between. \\
Quantum kernels & \cite{Havlicek2019,QuantumKernelDefect2022} & $\le 16$ & State-overlap kernel handed to a classical kernel method. \\
Symmetry-restricted circuits & \cite{Meyer2023,Nguyen2024EQNN,West2024,West2023} & varies & Generators/ansatz restricted to a symmetry group's commutant. \\
Hamiltonian-regularized hybrids & \cite{Shi2022,Park2024} & $\le 8$ & Physical energy functional as auxiliary training signal or primary object. \\
\bottomrule
\end{tabular}
\end{table}

\section{Open Challenges and Future Directions}
\label{sec:challenges}
Several open problems recur across the sections above rather than being specific to any one subfield. First, \emph{trainability at scale}: every architectural mitigation for barren plateaus surveyed in Section~\ref{sec:barren} has been demonstrated at qubit counts far below what would be needed to encode realistic image resolutions directly (Section~\ref{sec:qip}), so it remains open whether these mitigations compose favorably as circuit width grows toward that regime. Second, \emph{real-hardware validation} remains comparatively rare relative to simulator-only results across nearly every architecture surveyed here; where hardware demonstrations exist, they are typically small-scale feasibility pilots rather than statistically powered benchmarks with error mitigation and repeated trials. Third, \emph{standardized, resource-matched benchmarking} — comparing a quantum or hybrid model against a classical control matched in feature width and trainable parameter count, rather than an under-matched or task-mismatched classical baseline — remains inconsistently applied across the literature, despite being exactly the discipline the dequantization and power-of-data results of Section~\ref{sec:advantage} show is necessary to distinguish genuine advantage from architectural motivation. Fourth, \emph{integrating enforced (not merely motivated) symmetry constructions end-to-end into trainable pipelines} (Section~\ref{sec:symmetry}) remains, in most architectures surveyed, a post-hoc or diagnostic-level construction rather than a component of the trained model itself. Fifth, the QCNN (Section~\ref{sec:qcnn}) and tensor-network (Section~\ref{sec:tensor}) literatures have developed largely in parallel despite a close mathematical relationship — a QCNN's pooling structure and an MPS's canonical-form truncation are both, at heart, hierarchical coarse-graining procedures — and a more unified treatment connecting the two remains underdeveloped.

\section{Conclusion}
\label{sec:conclusion}
The literature surveyed here spans variational circuit design and its trainability obstructions, quantum and tensor-network representations of image data, quantum convolutional architectures, quantum autoencoders and dressed hybrid circuits, quantum kernel methods, symmetry and group-equivariance in classical and quantum form, physics-informed Hamiltonian regularization, and the quantum-advantage and dequantization results that bound what any of the preceding architectural choices can be expected to deliver. Read together, these strands support a specific methodological conclusion rather than a general one: architectural choices motivated by symmetry, physical structure, hierarchical coarse-graining, or entanglement are well-justified as design principles, and several (Sections~\ref{sec:qcnn}--\ref{sec:hamiltonian}) admit rigorous mathematical treatment establishing exactly what property they guarantee — equivariance, bounded gradient variance, compression fidelity, or exact solvability — but none of these guarantees is, by itself, a guarantee of empirical necessity or advantage on any specific downstream task, a distinction the dequantization and power-of-data literature (Section~\ref{sec:advantage}) makes precise and that resource-matched, leakage-audited empirical evaluation, rather than architectural motivation alone, is required to resolve.

\bibliographystyle{unsrt}
\begin{thebibliography}{99}

\bibitem{Cerezo2021} M. Cerezo et al., ``Variational quantum algorithms,'' \textit{Nature Reviews Physics}, vol. 3, no. 9, pp. 625--644, 2021. arXiv:2012.09265.
\bibitem{Peruzzo2014} A. Peruzzo, J. McClean, P. Shadbolt, et al., ``A variational eigenvalue solver on a photonic quantum processor,'' \textit{Nature Communications}, vol. 5, p. 4213, 2014. arXiv:1304.3061.
\bibitem{Tilly2022} J. Tilly et al., ``The Variational Quantum Eigensolver: A review of methods and best practices,'' \textit{Physics Reports}, vol. 986, pp. 1--128, 2022.
\bibitem{Bauer2020} B. Bauer, S. Bravyi, M. Motta, and G. K.-L. Chan, ``Quantum Algorithms for Quantum Chemistry and Quantum Materials Science,'' \textit{Chemical Reviews}, vol. 120, no. 22, pp. 12685--12717, 2020.
\bibitem{Bharti2022} K. Bharti, A. Cervera-Lierta, T. H. Kyaw, et al., ``Noisy intermediate-scale quantum (NISQ) algorithms,'' \textit{Reviews of Modern Physics}, vol. 94, p. 015004, 2022. arXiv:2101.08448.
\bibitem{Preskill2018} J. Preskill, ``Quantum Computing in the NISQ era and beyond,'' \textit{Quantum}, vol. 2, p. 79, 2018. arXiv:1801.00862.
\bibitem{McClean2018} J. R. McClean, S. Boixo, V. N. Smelyanskiy, R. Babbush, and H. Neven, ``Barren plateaus in quantum neural network training landscapes,'' \textit{Nature Communications}, vol. 9, no. 1, p. 4812, 2018. arXiv:1803.11173.
\bibitem{Pesah2021} A. Pesah, M. Cerezo, S. Wang, T. Volkoff, A. T. Sornborger, and P. J. Coles, ``Absence of Barren Plateaus in Quantum Convolutional Neural Networks,'' \textit{Physical Review X}, vol. 11, p. 041011, 2021. arXiv:2011.02966.
\bibitem{Larocca2025} M. Larocca, S. Thanasilp, S. Wang, et al., ``Barren Plateaus in Variational Quantum Computing,'' \textit{Nature Reviews Physics}, 2025. arXiv:2405.00781.
\bibitem{Park2024} C.-Y. Park and N. Killoran, ``Hamiltonian variational ansatz without barren plateaus,'' \textit{Quantum}, vol. 8, p. 1239, 2024. arXiv:2302.08529.
\bibitem{Cunningham2024} J. Cunningham and J. Zhuang, ``Investigating and Mitigating Barren Plateaus in Variational Quantum Circuits: A Survey,'' \textit{Quantum Information Processing}, 2024/2025. arXiv:2407.17706.
\bibitem{Deshpande2024} A. Deshpande, M. Hinsche, K. Najafi, K. Sharma, R. Sweke, and C. Zoufal, ``Dynamic parameterized quantum circuits: expressive and barren-plateau free,'' 2024 (rev.\ 2025). arXiv:2411.05760.
\bibitem{Zaman2023} K. Zaman, A. Marchisio, M. A. Hanif, and M. Shafique, ``A Survey on Quantum Machine Learning: Current Trends, Challenges, Opportunities, and the Road Ahead,'' 2023. arXiv:2310.10315.
\bibitem{Wang2024QML} Y. Wang and J. Liu, ``A comprehensive review of Quantum Machine Learning: from NISQ to Fault Tolerance,'' \textit{Reports on Progress in Physics}, vol. 87, p. 116402, 2024. arXiv:2401.11351.
\bibitem{Le2011} P. Q. Le, F. Dong, and K. Hirota, ``A flexible representation of quantum images for polynomial preparation, image compression, and processing operations,'' \textit{Quantum Information Processing}, vol. 10, pp. 63--84, 2011.
\bibitem{Zhang2013} Y. Zhang, K. Lu, Y. Gao, and M. Wang, ``NEQR: A novel enhanced quantum representation of digital images,'' \textit{Quantum Information Processing}, vol. 12, no. 8, pp. 2833--2860, 2013.
\bibitem{Zhou2017} R. Zhou, W. Hu, P. Fan, and H. Ian, ``Quantum realization of the bilinear interpolation method for NEQR,'' \textit{Scientific Reports}, vol. 7, no. 1, p. 2511, 2017.
\bibitem{Yao2017} X. Yao et al., ``Quantum Image Processing and Its Application to Edge Detection: Theory and Experiment,'' \textit{Physical Review X}, vol. 7, no. 3, p. 031041, 2017.
\bibitem{QIPSystematicReview2025} ``A systematic review of quantum image processing: Representation, applications and future perspectives,'' \textit{Information Fusion}, 2025.
\bibitem{QuantumImageRepClassification2025} ``Analysis of Quantum Image Representations for Supervised Classification,'' 2025. arXiv:2507.22039.
\bibitem{Chakraborty2018} S. Chakraborty, S. B. Mandal, and S. H. Shaikh, ``Quantum image processing: challenges and future research issues,'' \textit{International Journal of Information Technology}, vol. 14, no. 1, pp. 475--489, 2018.
\bibitem{Cong2019} I. Cong, S. Choi, and M. D. Lukin, ``Quantum Convolutional Neural Networks,'' \textit{Nature Physics}, vol. 15, pp. 1273--1278, 2019. arXiv:1810.03787.
\bibitem{Hur2022} T. Hur, L. Kim, and D. K. Park, ``Quantum convolutional neural network for classical data classification,'' \textit{Quantum Machine Intelligence}, vol. 4, art. 3, 2022. arXiv:2108.00661.
\bibitem{Cirac2021} J. I. Cirac, D. P\'erez-Garc\'ia, N. Schuch, and F. Verstraete, ``Matrix product states and projected entangled pair states: Concepts, symmetries, theorems,'' \textit{Reviews of Modern Physics}, vol. 93, no. 4, p. 045003, 2021.
\bibitem{Latorre2005} J. I. Latorre, ``Image compression and entanglement,'' 2005. arXiv:quant-ph/0510031.
\bibitem{Liu2021} Y. Liu, W.-J. Li, X. Zhang, M. Lewenstein, G. Su, and S.-J. Ran, ``Entanglement-Based Feature Extraction by Tensor Network Machine Learning,'' \textit{Frontiers in Applied Mathematics and Statistics}, 2021.
\bibitem{Jobst2024} B. Jobst, K. Shen, C. A. Riofr\'io, E. Shishenina, and F. Pollmann, ``Efficient MPS representations and quantum circuits from the Fourier modes of classical image data,'' \textit{Quantum}, vol. 8, p. 1544, 2024. arXiv:2311.07666.
\bibitem{Stoudenmire2016} E. M. Stoudenmire and D. J. Schwab, ``Supervised Learning with Tensor Networks,'' in \textit{Advances in Neural Information Processing Systems}, vol. 29, 2016. arXiv:1605.05775.
\bibitem{Han2018} Z. Han, J. Wang, H. Fan, L. Wang, and P. Zhang, ``Unsupervised Generative Modeling Using Matrix Product States,'' \textit{Physical Review X}, vol. 8, no. 3, p. 031012, 2018. arXiv:1709.01662.
\bibitem{Guala2023} D. Guala, S. Zhang, E. Cervero-Mart\'in, K. Chumak, C. Blank, and C. Ke, ``A Practical Overview of Image Classification with Variational Tensor-Network Quantum Circuits,'' \textit{Scientific Reports}, vol. 13, no. 1, p. 4427, 2023. arXiv:2209.11058.
\bibitem{Wang2023TN} M. Wang, Y. Pan, Z. Xu, G. Li, X. Yang, D. Mandic, and A. Cichocki, ``Tensor Networks Meet Neural Networks: A Survey and Future Perspectives,'' 2023. arXiv:2302.09019.
\bibitem{Araz2022} J. Y. Araz and M. Spannowsky, ``Classical versus quantum: Comparing tensor-network-based quantum circuits on Large Hadron Collider data,'' \textit{Physical Review A}, vol. 106, no. 6, p. 062423, 2022.
\bibitem{Romero2017} J. Romero, J. P. Olson, and A. Aspuru-Guzik, ``Quantum autoencoders for efficient compression of quantum data,'' \textit{Quantum Science and Technology}, vol. 2, p. 045001, 2017. arXiv:1612.02806.
\bibitem{WangQAECompression2024} ``Quantum image compression with autoencoders based on parameterized quantum circuits,'' \textit{Quantum Information Processing}, 2024.
\bibitem{QAEClassification2025} ``Quantum autoencoders for image classification,'' 2025. arXiv:2502.15254.
\bibitem{QCAE2023} ``Quantum Circuit AutoEncoder,'' 2023. arXiv:2307.08446.
\bibitem{Mari2020} A. Mari, T. R. Bromley, J. Izaac, M. Schuld, and N. Killoran, ``Transfer learning in hybrid classical-quantum neural networks,'' \textit{Quantum}, vol. 4, p. 340, 2020. arXiv:1912.08278.
\bibitem{Zoufal2019} C. Zoufal, A. Lucchi, and S. Woerner, ``Quantum Generative Adversarial Networks for Learning and Loading Random Distributions,'' \textit{npj Quantum Information}, vol. 5, p. 103, 2019. arXiv:1904.00043.
\bibitem{HuangExpGAN2020} H.-L. Huang et al., ``Experimental Quantum Generative Adversarial Networks for Image Generation,'' 2020. arXiv:2010.06201.
\bibitem{HybridGANHighRes2022} ``Hybrid Quantum-Classical Generative Adversarial Network for High Resolution Image Generation,'' 2022. arXiv:2212.11614.
\bibitem{MosaiQ2023} ``MosaiQ: Quantum Generative Adversarial Networks for Image Generation on NISQ Computers,'' 2023. arXiv:2308.11096.
\bibitem{QuantumDiffusion2024} ``Quantum Denoising Diffusion Models,'' 2024. arXiv:2401.07049.
\bibitem{HybridLatentDiffusion2025} ``Hybrid Quantum-Classical Latent Diffusion Models for Medical Image Generation,'' 2025. arXiv:2508.09903.
\bibitem{Havlicek2019} V. Havl{\'i}{\v c}ek, A. D. C{\'o}rcoles, K. Temme, A. W. Harrow, A. Kandala, J. M. Chow, and J. M. Gambetta, ``Supervised learning with quantum-enhanced feature spaces,'' \textit{Nature}, vol. 567, no. 7747, pp. 209--212, 2019. arXiv:1804.11326.
\bibitem{QuantumKernelDefect2022} ``Quantum Kernel for Image Classification of Real World Manufacturing Defects,'' 2022. arXiv:2212.08693.
\bibitem{Slabbert2024} D. Slabbert and F. Petruccione, ``Hybrid Quantum-Classical Feature Extraction approach for Image Classification using Autoencoders and Quantum SVMs,'' 2024. arXiv:2410.18814.
\bibitem{Cohen2016} T. Cohen and M. Welling, ``Group Equivariant Convolutional Networks,'' in \textit{Proceedings of the 33rd International Conference on Machine Learning}, PMLR, vol. 48, 2016, pp. 2990--2999. arXiv:1602.07576.
\bibitem{Weiler2019} M. Weiler and G. Cesa, ``General E(2)-Equivariant Steerable CNNs,'' in \textit{Advances in Neural Information Processing Systems}, vol. 32, 2019. arXiv:1911.08251.
\bibitem{Bronstein2021} M. M. Bronstein, J. Bruna, T. Cohen, and P. Veli\v{c}kovi\'c, ``Geometric Deep Learning: Grids, Groups, Graphs, Geodesics, and Gauges,'' 2021. arXiv:2104.13478.
\bibitem{Meyer2023} J. J. Meyer, M. Mularski, E. Gil-Fuster, A. A. Mele, F. Arzani, A. Wilms, and J. Eisert, ``Exploiting Symmetry in Variational Quantum Machine Learning,'' \textit{PRX Quantum}, vol. 4, p. 010328, 2023. arXiv:2205.06217.
\bibitem{Nguyen2024EQNN} Q. T. Nguyen, L. Schatzki, P. Braccia, M. Ragone, P. J. Coles, F. Sauvage, M. Larocca, and M. Cerezo, ``Theory for Equivariant Quantum Neural Networks,'' \textit{PRX Quantum}, vol. 5, p. 020328, 2024. arXiv:2210.08566.
\bibitem{Schatzki2024} L. Schatzki, M. Larocca, F. Sauvage, and M. Cerezo, ``Theoretical Guarantees for Permutation-Equivariant Quantum Neural Networks,'' \textit{npj Quantum Information}, vol. 10, p. 12, 2024. arXiv:2210.09974.
\bibitem{Skolik2023} A. Skolik, M. Cattelan, S. Yarkoni, T. B\"ack, and V. Dunjko, ``Equivariant Quantum Circuits for Learning on Weighted Graphs,'' \textit{npj Quantum Information}, vol. 9, p. 47, 2023. arXiv:2205.06109.
\bibitem{Sauvage2024} F. Sauvage, M. Larocca, P. J. Coles, and M. Cerezo, ``Building Spatial Symmetries into Parameterized Quantum Circuits for Faster Training,'' \textit{Quantum Science and Technology}, vol. 9, p. 015029, 2024. arXiv:2207.14413.
\bibitem{West2024} M. West, J. Heredge, M. Sevior, and M. Usman, ``Provably Trainable Rotationally Equivariant Quantum Machine Learning,'' \textit{PRX Quantum}, vol. 5, p. 030320, 2024. arXiv:2311.05873.
\bibitem{West2023} M. West, M. Sevior, and M. Usman, ``Reflection Equivariant Quantum Neural Networks for Enhanced Image Classification,'' \textit{Machine Learning: Science and Technology}, vol. 4, p. 035027, 2023. arXiv:2212.00264.
\bibitem{West2025} M. West, M. Sevior, and M. Usman, ``Image Classification with Rotation-Invariant Variational Quantum Circuits,'' \textit{Physical Review Research}, vol. 7, p. 013082, 2025. arXiv:2403.15031.
\bibitem{Craps2024} B. Craps, M. D. Clerck, O. Evnin, and P. Hacker, ``Integrability and complexity in quantum spin chains,'' \textit{SciPost Physics}, vol. 16, no. 2, p. 041, 2024.
\bibitem{Franchini2017} F. Franchini, \textit{An Introduction to Integrable Techniques for One-Dimensional Quantum Systems}, Lecture Notes in Physics, vol. 940. Springer, 2017. arXiv:1609.02100.
\bibitem{Shi2022} J. Shi, W. Wang, X. Lou, S. Zhang, and X. Li, ``Parameterized Hamiltonian Learning With Quantum Circuit,'' \textit{IEEE Transactions on Pattern Analysis and Machine Intelligence}, vol. 45, no. 5, pp. 6086--6095, 2022.
\bibitem{Chinni2023} G. Chinni et al., ``Physics Informed Neural Networks Learning a Two-Qubit Hamiltonian,'' 2023. arXiv:2310.15148.
\bibitem{Tang2019} E. Tang, ``A quantum-inspired classical algorithm for recommendation systems,'' in \textit{Proceedings of the 51st Annual ACM SIGACT Symposium on Theory of Computing (STOC)}, 2019. arXiv:1807.04271.
\bibitem{Cotler2021} J. Cotler, H.-Y. Huang, and J. R. McClean, ``Revisiting dequantization and quantum advantage in learning tasks,'' 2021. arXiv:2112.00811.
\bibitem{Gyurik2023} C. Gyurik and V. Dunjko, ``Exponential separations between classical and quantum learners,'' 2023. arXiv:2306.16028.
\bibitem{Huang2021PowerData} H.-Y. Huang, M. Broughton, M. Mohseni, R. Babbush, S. Boixo, H. Neven, and J. R. McClean, ``Power of data in quantum machine learning,'' \textit{Nature Communications}, vol. 12, no. 1, p. 2631, 2021. arXiv:2011.01938.
\bibitem{Huang2020Shadows} H.-Y. Huang, R. Kueng, and J. Preskill, ``Predicting many properties of a quantum system from very few measurements,'' \textit{Nature Physics}, vol. 16, pp. 1050--1057, 2020.

\end{thebibliography}

\end{document}
